\documentclass[12pt,a4paper]{article}
\usepackage[a4paper,left=23mm,right=20mm,top=24mm,bottom=24mm]{geometry}
\usepackage[T2A]{fontenc}
\usepackage[utf8]{inputenc}
\usepackage[russian]{babel}
\usepackage{amsmath,amssymb,mathtools}
\usepackage{array,booktabs,longtable}
\usepackage{xcolor}
\usepackage{tikz}
\usepackage[unicode,colorlinks=true,linkcolor=blue!55!black]{hyperref}
\setlength{\parindent}{0pt}
\setlength{\parskip}{0.32em}
\setlength{\emergencystretch}{3em}
\definecolor{muted}{RGB}{72,78,86}
\definecolor{accent}{RGB}{35,74,99}
\newcommand{\reviewfield}[1]{\par\smallskip\noindent\textbf{#1}\par}
\newcommand{\mathunit}[1]{\,\text{#1}}
\hypersetup{pdftitle={Ревью домашней работы},pdfauthor={Лёвин Артём Александрович}}

\begin{document}
\thispagestyle{empty}
\vspace*{0.26\textheight}
\begin{center}
{\Huge\bfseries Ревью домашней работы\par}
\vspace{2.2cm}
{\large Лёвин Артём Александрович\par}
\vspace{0.65cm}
{\large 25 октября 2026 года\par}
\end{center}
\vfill
\begin{center}
\begin{tikzpicture}
\draw[line width=0.8pt,draw=accent!55] (-4,0) .. controls (-2,0.65) and (1,-0.60) .. (4,0);
\end{tikzpicture}
\end{center}
\newpage
\section*{Сводная таблица}
\small
\renewcommand{\arraystretch}{1.4}
\begin{tabular}{@{}>{\raggedright\arraybackslash}p{0.09\linewidth}>{\raggedright\arraybackslash}p{0.18\linewidth}>{\raggedright\arraybackslash}p{0.25\linewidth}>{\raggedright\arraybackslash}p{0.16\linewidth}>{\raggedright\arraybackslash}p{0.18\linewidth}@{}}
\toprule
\textbf{№} & \textbf{Тема} & \textbf{Суть ошибки} & \textbf{Ваш ответ} & \textbf{Правильный ответ} \\
\midrule
\hyperref[task:1]{1} & Второй закон Ньютона в проекциях & Повторно учтена сила тяжести; добавлено лишнее слагаемое & $-7\,\text{м/с}^{2}$ & $3\,\text{м/с}^{2}$, вверх \\
\addlinespace
\hyperref[task:9]{9} & Закон Архимеда и равномерное движение & Решение не представлено & Не указан & $0{,}192\,\text{м}^{3}$ \\
\addlinespace
\hyperref[task:10]{10} & Второй закон Ньютона и равноускоренное движение & Решение не представлено & Не указан & $8\,\text{м/с}^{2}$, вниз; $2\mathunit{с}$ \\
\bottomrule
\end{tabular}
\normalsize
\vfill
\noindent{\small\color{muted}В таблице и подробном разборе приведены только задания с ошибками или отсутствующими решениями.}
\newpage
\thispagestyle{empty}
\null\vfill
\begin{center}
{\Huge\bfseries Разбор ошибок}
\end{center}
\vfill\null
\newpage

\phantomsection\label{task:1}
\section*{Задание № 1. Проекции второго закона Ньютона}
\reviewfield{Текст задания}
На тело массой $2\mathunit{кг}$ действуют сила $26\mathunit{Н}$ вверх и сила $20\mathunit{Н}$ вниз. Ось $Oy$ направлена вверх. Запишите второй закон Ньютона в проекциях, определите $a_y$ и направление ускорения.

\reviewfield{Ваше решение}
В записи правильно задана вертикальная ось, однако при составлении уравнения к двум названным силам прибавлена ещё одна сила тяжести:
\[
ma_y=F_1-F_2-mg.
\]
Затем подставлены числа:
\[
a_y=\frac{26-20-2\cdot10}{2}=-7\,\text{м/с}^{2}.
\]

\reviewfield{Ваш ответ}
$a_y=-7\,\text{м/с}^{2}$.

\reviewfield{Ошибка}
\textcolor{red}{Ошибка:} в уравнение добавлено слагаемое $-mg$, которого нет в перечне двух действующих сил как отдельной, третьей силы.

\reviewfield{Где именно возникает ошибка}
\textbf{Последний правильный шаг:} ось $Oy$ направлена вверх, поэтому проекция силы $26\mathunit{Н}$ положительна, а проекция силы $20\mathunit{Н}$ отрицательна.

\textbf{Первый неправильный шаг:} вместо суммы проекций \(F_1-F_2\) записано
\[
\underbrace{F_1-F_2}_{\text{две силы из условия}}\;\color{red}{-mg}.
\]
В результате вниз направленная сила была фактически учтена ещё раз.

\reviewfield{Почему это неверно}
Второй закон Ньютона требует сложить \emph{проекции всех сил, которые реально действуют на тело}: сумма проекций равна произведению массы на проекцию ускорения.

Здесь условие сообщает о двух силах, численно равных $26\mathunit{Н}$ и $20\mathunit{Н}$, и задаёт их направления. Если сила $20\mathunit{Н}$ представляет собой вес тела, она уже учтена в условии; если она имеет другую физическую природу, добавлять третью силу без соответствующего указания также нельзя. Поэтому самостоятельное добавление $mg$ не обосновано.

Из-за лишнего отрицательного слагаемого суммарная сила получилась направленной вниз, а ускорение --- отрицательным относительно выбранной оси. Это меняет не только величину, но и физический вывод о направлении ускорения.

\reviewfield{Ключевая идея}
\textcolor{blue!60!black}{Ключевая идея:} сначала перечислите все силы из условия ровно по одному разу, затем поставьте знак каждой проекции с учётом направления оси и только после этого применяйте второй закон Ньютона.

\reviewfield{Исправление от последнего правильного шага}
При сохранении верно выбранной оси $Oy$ и знаков двух сил получаем
\[
ma_y=F_1-F_2.
\]
Вычисляем результирующую силу и проекцию ускорения:
\[
F_{y,\text{рез}}=26-20=6\mathunit{Н},
\qquad
 a_y=\frac{6}{2}=3\,\text{м/с}^{2}.
\]

\reviewfield{Полное правильное решение}
\textbf{Шаг 1.} Выбираем положительное направление оси $Oy$ вверх, как требует условие.

\textbf{Шаг 2.} Выписываем проекции всех названных сил: первая сила направлена вверх и имеет проекцию $+26\mathunit{Н}$; вторая направлена вниз и имеет проекцию $-20\mathunit{Н}$.

\textbf{Шаг 3.} Записываем второй закон Ньютона в проекции на ось $Oy$:
\[
ma_y=F_{1y}+F_{2y}=26-20.
\]
\textbf{Шаг 4.} Подставляем массу тела и находим ускорение:
\[
2a_y=6,
\qquad a_y=\frac{6}{2}=3\,\text{м/с}^{2}.
\]
\textbf{Шаг 5.} Проекция положительна. Следовательно, ускорение направлено \emph{вверх}, то есть в положительную сторону оси $Oy$.

\reviewfield{Проверка}
По найденному ускорению произведение массы на ускорение равно
\[
ma_y=2\cdot3=6\mathunit{Н}.
\]
Сумма проекций данных сил также равна
\[
26-20=6\mathunit{Н}.
\]
Левая и правая части второго закона Ньютона совпадают; направление суммарной силы соответствует направлению ускорения.

\reviewfield{Итог}
\textcolor{teal!60!black}{Итог:} проекция ускорения $a_y=3\,\text{м/с}^{2}$; ускорение направлено вверх.

\reviewfield{Задача на закрепление}
На тело массой $2\mathunit{кг}$ действуют только две вертикальные силы: $23\mathunit{Н}$ вверх и $17\mathunit{Н}$ вниз. Ось $Oy$ направлена вверх. Запишите второй закон Ньютона в проекциях, определите проекцию ускорения и его направление.
\newpage

\phantomsection\label{task:9}
\section*{Задание № 9. Закон Архимеда и равновесие сил}
\reviewfield{Текст задания}
Воздушный шар с суммарной массой $0{,}2\mathunit{кг}$ поднимается равномерно в неподвижном воздухе. Сопротивление равно $0{,}4\mathunit{Н}$. При плотности воздуха $\rho_{\text{возд}}=1{,}25\,\text{кг/м}^{3}$ найдите объём вытесненного воздуха.

\reviewfield{Ваше решение}
На предоставленных фотографиях решение этого задания отсутствует.

\reviewfield{Ваш ответ}
Не указан.

\reviewfield{Ошибка}
\textcolor{red}{Ошибка:} задание осталось без решения и численного ответа. Здесь не обнаружен ошибочный математический переход: проверять последовательность вычислений невозможно, поскольку она не представлена.

\reviewfield{Где именно возникает неполнота}
\textbf{Последний достоверно выполненный шаг:} на фотографиях не зафиксирован.

\textbf{Первый отсутствующий необходимый шаг:} запись условия равномерного подъёма через второй закон Ньютона:
\[
F_{\text{А}}-mg-F_{\text{сопр}}=0.
\]
Именно это уравнение связывает искомый объём с известными величинами.

\reviewfield{Почему без этого шага нельзя получить ответ}
При равномерном подъёме скорость шара постоянна, а ускорение равно нулю. Это означает, что сумма всех сил равна нулю, хотя сами силы не исчезают.

На шар действуют три силы: сила Архимеда направлена вверх; сила тяжести направлена вниз; сила сопротивления воздуха направлена против движения, то есть тоже вниз. Силу Архимеда определяют через плотность окружающего воздуха и объём вытесненного воздуха:
\[
F_{\text{А}}=\rho_{\text{возд}}gV.
\]
Недостаточно приравнять силу Архимеда только весу шара: дополнительно требуется компенсировать сопротивление $0{,}4\mathunit{Н}$. Поэтому в данном случае $F_{\text{А}}=mg+F_{\text{сопр}}$, а не просто $mg$.

\reviewfield{Ключевая идея}
\textcolor{blue!60!black}{Ключевая идея:} равномерное движение означает нулевое ускорение; направленная вверх сила Архимеда уравновешивает вместе взятые силу тяжести и силу сопротивления.

\reviewfield{Исправление}
Поскольку собственного решения нет, исправление начинаем с необходимого уравнения равновесия сил. Выбираем направление вверх положительным, записываем $F_{\text{А}}-mg-F_{\text{сопр}}=0$ и заменяем силу Архимеда произведением $\rho_{\text{возд}}gV$.

\reviewfield{Полное правильное решение}
\textbf{Шаг 1.} Выписываем данные: масса шара $m=0{,}2\mathunit{кг}$, сила сопротивления $F_{\text{сопр}}=0{,}4\mathunit{Н}$, плотность воздуха $\rho_{\text{возд}}=1{,}25\,\text{кг/м}^{3}$. Примем школьное приближение $g=10\,\text{м/с}^{2}$.

\textbf{Шаг 2.} Так как шар поднимается \emph{равномерно}, его ускорение нулевое: $a_y=0$. Выбираем ось $Oy$ вверх.

\textbf{Шаг 3.} Проецируем силы на вертикальную ось:
\[
ma_y=F_{\text{А}}-mg-F_{\text{сопр}}.
\]
Из $a_y=0$ следует
\[
F_{\text{А}}=mg+F_{\text{сопр}}.
\]
\textbf{Шаг 4.} Используем закон Архимеда:
\[
\rho_{\text{возд}}gV=mg+F_{\text{сопр}}.
\]
Выражаем объём:
\[
V=\frac{mg+F_{\text{сопр}}}{\rho_{\text{возд}}g}.
\]
\textbf{Шаг 5.} Находим вес шара:
\[
mg=0{,}2\cdot10=2\mathunit{Н}.
\]
Вычисляем требуемую силу Архимеда:
\[
F_{\text{А}}=2+0{,}4=2{,}4\mathunit{Н}.
\]
Подставляем значения в формулу объёма:
\[
V=\frac{2{,}4}{1{,}25\cdot10}
 =\frac{2{,}4}{12{,}5}
 =0{,}192\,\text{м}^{3}.
\]

\reviewfield{Проверка}
Проверим, какую выталкивающую силу даёт найденный объём:
\[
F_{\text{А}}=1{,}25\cdot10\cdot0{,}192=2{,}4\mathunit{Н}.
\]
Сумма сил, направленных вниз, равна $mg+F_{\text{сопр}}=2+0{,}4=2{,}4\mathunit{Н}$. Следовательно, равнодействующая равна нулю, что согласуется с равномерным подъёмом. Единицы измерения результата --- кубические метры, как и требуется для объёма.

\reviewfield{Итог}
\textcolor{teal!60!black}{Итог:} объём вытесненного воздуха равен $0{,}192\,\text{м}^{3}$.

\reviewfield{Задача на закрепление}
Воздушный шар общей массой $0{,}3\mathunit{кг}$ равномерно поднимается в неподвижном воздухе. Сила сопротивления воздуха равна $0{,}6\mathunit{Н}$. Плотность воздуха составляет $1{,}2\,\text{кг/м}^{3}$. Определите объём вытесненного воздуха, принимая $g=10\,\text{м/с}^{2}$.
\newpage

\phantomsection\label{task:10}
\section*{Задание № 10. Ускорение и время падения}
\reviewfield{Текст задания}
Тело массой $0{,}25\mathunit{кг}$ отпускают без начальной скорости с высоты $16\mathunit{м}$. Постоянная сила сопротивления равна $0{,}5\mathunit{Н}$ и направлена вверх. Найдите ускорение и время падения. Силой Архимеда пренебречь.

\reviewfield{Ваше решение}
На предоставленных фотографиях решение этого задания отсутствует.

\reviewfield{Ваш ответ}
Не указан.

\reviewfield{Ошибка}
\textcolor{red}{Ошибка:} не выполнены два обязательных пункта задания: нахождение ускорения с учётом сопротивления и вычисление времени падения. Ошибочные промежуточные записи отсутствуют, потому что само решение не представлено.

\reviewfield{Где именно возникает неполнота}
\textbf{Последний достоверно выполненный шаг:} на фотографиях не зафиксирован.

\textbf{Первый отсутствующий необходимый шаг:} проекция второго закона Ньютона при положительном направлении оси вниз:
\[
ma=mg-F_{\text{сопр}}.
\]
После определения ускорения требуется применить формулу пути для движения без начальной скорости.

\reviewfield{Почему это принципиально}
Тело движется вниз, но сила сопротивления направлена вверх, то есть против движения. Поэтому сопротивление уменьшает равнодействующую силу и ускорение по сравнению со свободным падением. Если ошибочно сложить силы, ускорение окажется завышено; если принять ускорение равным $g$, сопротивление останется без учёта.

Здесь сила сопротивления постоянна, а масса и ускорение свободного падения принимаются постоянными. Поэтому равнодействующая и ускорение также постоянны; можно использовать школьные формулы равноускоренного движения. Поскольку тело \emph{отпускают} без начальной скорости, начальная скорость равна нулю.

\reviewfield{Ключевая идея}
\textcolor{blue!60!black}{Ключевая идея:} сначала по второму закону Ньютона найдите постоянное ускорение вниз с учётом противодействующей силы, затем подставьте это ускорение в закон движения тела из состояния покоя.

\reviewfield{Исправление}
Принимаем положительное направление оси вниз. Вначале находим равнодействующую $mg-F_{\text{сопр}}$, делим её на массу, а затем из равенства $h=at^2/2$ определяем положительное время падения.

\reviewfield{Полное правильное решение}
\textbf{Шаг 1.} Записываем данные:
\[
m=0{,}25\mathunit{кг},\quad h=16\mathunit{м},\quad
F_{\text{сопр}}=0{,}5\mathunit{Н},\quad v_0=0.
\]
Примем $g=10\,\text{м/с}^{2}$. Выбираем направление вниз положительным.

\textbf{Шаг 2.} На тело действуют сила тяжести $mg$ вниз и сила сопротивления $F_{\text{сопр}}$ вверх. По второму закону Ньютона:
\[
ma=mg-F_{\text{сопр}}.
\]
\textbf{Шаг 3.} Находим силу тяжести:
\[
mg=0{,}25\cdot10=2{,}5\mathunit{Н}.
\]
Равнодействующая направлена вниз и равна
\[
F_{\text{рез}}=2{,}5-0{,}5=2\mathunit{Н}.
\]
\textbf{Шаг 4.} Вычисляем ускорение:
\[
a=\frac{F_{\text{рез}}}{m}=\frac{2}{0{,}25}=8\,\text{м/с}^{2}.
\]
Оно направлено вниз, так как результат положителен для выбранной оси.

\textbf{Шаг 5.} Тело начинает движение из состояния покоя, поэтому при постоянном ускорении пройденный путь равен
\[
h=v_0t+\frac{at^2}{2}=\frac{at^2}{2}.
\]
Подставляем высоту и ускорение:
\[
16=\frac{8t^2}{2}=4t^2,
\qquad t^2=4.
\]
Время не может быть отрицательным, поэтому
\[
t=\sqrt{4}=2\mathunit{с}.
\]

\reviewfield{Проверка}
Проверяем второй закон Ньютона: $ma=0{,}25\cdot8=2\mathunit{Н}$, а разность силы тяжести и сопротивления равна $2{,}5-0{,}5=2\mathunit{Н}$.

Проверяем пройденный путь за найденное время:
\[
\frac{at^2}{2}=\frac{8\cdot2^2}{2}=16\mathunit{м},
\]
что совпадает с данной высотой. Ускорение $8\,\text{м/с}^{2}$ меньше $g=10\,\text{м/с}^{2}$, как и должно быть при сопротивлении, направленном вверх.

\reviewfield{Итог}
\textcolor{teal!60!black}{Итог:} ускорение тела равно $8\,\text{м/с}^{2}$ и направлено вниз; время падения равно $2\mathunit{с}$.

\reviewfield{Задача на закрепление}
Тело массой $0{,}5\mathunit{кг}$ отпускают без начальной скорости с высоты $36\mathunit{м}$. Во время падения на него действует постоянная сила сопротивления $1\mathunit{Н}$, направленная вверх. Силой Архимеда пренебречь. Найдите ускорение и время падения, принимая $g=10\,\text{м/с}^{2}$.

\vfill
\begin{center}
{\small\color{muted}Лёвин Артём Александрович эксклюзивно для Володи Хачатуряна}
\end{center}
\end{document}